\section{Conclusion}
In this chapter, we have designed a simple yet efficient solution to the problem of constant delay in the observation of the state or in the execution of the action.
This solution consists in first learning a policy in the undelayed environment before imitating this policy in the delayed environment.
We proved theoretical bounds on the performance of such an imitated policy with respect to the undelayed expert when the environment demonstrates smoothness. 
Using these theoretical insights, we designed DIDA, an algorithm that follows the above principle, using \textsc{DAgger} as its imitation algorithm. 
The experimental results demonstrate the efficiency of DIDA compared to many baselines.
DIDA almost always achieves superior returns while requiring fewer samples.
Interestingly, we have shown how this algorithm can be modified to be used in three more scenarios.
First, when the delay is constant but not an integer, we provide a simple adjustment to DIDA that shows good empirical results. 
Second, while DIDA is trained on constant delays, we have provided theoretical guarantees when tested on a stochastic delay instead.
Third, we have shown that the simple ``trick'' of \Cref{chap:belief_based} can be applied to DIDA to leverage a policy learnt on a certain delay and efficiently apply it on smaller delays. 
As a future research direction, the current algorithm could be modified to be trained on stochastic delays and learn the policy of \Cref{eq:belief_pol_stoch}. 
In this way, it would obtain better theoretical guarantees than the one provided in \Cref{thm:stoch_mdp_dida_bound}.
Another direction would be to tackle DIDA's main drawback, the necessity of the undelayed expert. 
A potential idea could be to learn an undelayed expert offline from a dataset collected by a delayed policy. 
This would imply that no sample is ever needed from an undelayed environment.